A RAG system is best read as an evidence-conditioned decision process: a query arrives, a
retriever assembles a context, and a generator commits to an action on the basis of that
context. Failures occur when the assembled evidence is incomplete, irrelevant, misleading,
mutually contradictory, corrupted, or simply insufficient to support any answer
(Appendix~\ref{app:background}). Sorting the literature by \emph{where} in this process a
paper intervenes yields four failure surfaces, which we use as the top-level organisation of
the survey and summarise in Figure~\ref{fig:survey_outline}.

\textbf{Evidence-centric} work (\S\ref{sec:evidence}) treats the composition of the context
as the object of study, asking which passages reach the generator and in what order.
\textbf{Model-centric} work (\S\ref{sec:model}) holds the context fixed and asks how the
generator uses it, which is where faithfulness, attribution and knowledge conflict live.
\textbf{Adversarial} work (\S\ref{sec:adversarial}) replaces the assumption of benign noise
with an adaptive attacker who chooses the evidence. \textbf{System-centric} work
(\S\ref{sec:system}) steps outside the single forward pass and asks what the system should
\emph{do}---retrieve again, verify, abstain, or answer---under uncertainty and cost.

The four surfaces partition the literature, not the phenomena, and a single event routinely
touches several at once. A poisoned passage is an evidence state, an adversary capability,
and---if the system declines to answer because of it---a policy decision. That overlap is
exactly why placing two papers on the same surface does not make their numbers comparable,
and it motivates the reduction to five reporting coordinates in \S\ref{sec:coordinates}.
Evaluation practice cuts across all four surfaces and is treated separately in
Appendix~\ref{app:evaluation}; the rules used to assign a paper to a surface are in
Appendix~\ref{sec:appendix-methods}.
